\documentclass{ctexbook}
\usepackage{amsmath}
\usepackage{bm}
\usepackage{amssymb}
\title{高等数学上期末Review}
\begin{document}
	\part{一元函数微分学}
	\chapter{极限}
	\section{前置知识：集合与映射}
	\subsection{集合}
	\paragraph{集合的定义}~{}
	\\
	\indent
	我们在讨论问题时，在一定范围内说的对象，如数、向量、多项式等，笼统的称为\textbf{元素}或\textbf{元}.
	若干个固定元素的全体，叫做一个\textbf{集合}，简称为\textbf{集}.
	
	\paragraph{集合的分类}~{}
	\\
	\indent
	按照集合内元素是否有限，可将集合分为\textbf{有限集}和\textbf{无限集}.
	
	集合$A$内元素的个数用$|A|$表示，显然若$A$为有限集，$|A|<\infty$，反之，$A$则为无限集.
	
	若集合$A$满足$|A|=0$，则称集合$A$为\textbf{空集}，记为$\varnothing$.
	
	我们给出元素为数的几个常用集合.自然数集$\mathbb{N}$，正整数集$\mathbb{N^+}$，整数集$\mathbb{Z}$，有理数集$\mathbb{Q}$，非零有理数集$\mathbb{Q^*}$，实数集$\mathbb{R}$，非零实数集$\mathbb{R^*}$，复数集$\mathbb{C}$，非零复数集$\mathbb{C^*}$.
	
	\paragraph{集合的表述}~{}
	\\
	\indent
	集合的描述总体有\textbf{描述法}和\textbf{列举法}（或枚举法）.
	
	枚举法比较直接，直接将集合中的元素写出来，常用于有限集和有规律的无限集.例如四元数集$G=\{1,\mathrm{i},\mathrm{j},\mathrm{k},-1,-\mathrm{i},-\mathrm{j},-\mathrm{k}\}$.
	
	描述法则采用集合中元素公共属性的方式来描述集合，即：
	
	\centerline{$S=\{x|x$具有性质$P\}$}
	
	例如有理数集$\mathbb{Q}$可用描述法表示为
	
	\centerline{$\mathbb{Q}=\{x|x=\dfrac{p}{q}$,其中$p\in\mathbb{N^*}$并且$q\in\mathbb{Z}\}$}.
	
	\paragraph{两集合之间的关系}~{}
	\\
	\indent
	设两个集合$S,T$，若$S$中的元素均属于$T$，即
	
	\centerline{$x\in S \Rightarrow x\in T$}
	\noindent{则称$S$是$T$的\textbf{子集}，记为$S\subset T$.}
	
	若存在$x\in T$但$x\notin S$（即$T \not\subset S$），且满足$S\subset T$，则称$S$是$T$的\textbf{真子集}，记为$S \subsetneqq T$.
	
	同时，我们定义两集合相等，当且仅当两集合互为对方集合的子集，即
	
	\centerline{$S=T \Leftrightarrow S\subset T$且$T \subset S$.}
	
	\paragraph{集合的运算}~{}
	\\
	\indent
	由集合$S$和集合$T$的公共元素做成的集合，叫做$S$与$T$的\textbf{交集}，记为$S\cup T$.
	
	由集合$S$和集合$T$的所有元素做成的集合，叫做$S$与$T$的\textbf{并集}，记为$S\cap T$.
	
	集合的交并运算有一些基本的性质：
	
	（1）反身律：$A\cup A=A$,$\,A\cap A=A$;
	
	（2）交换律：$A\cup B=B\cup A$,$\,A\cap B=B\cap A$;
	
	（3）结合律：$A\cup (B \cup C)=(A\cup B)\cup C$,$\,A\cap (B \cap C)=(A\cap B)\cap C$;
	
	（4）分配律：
	\noindent{
	\begin{align}
		A\cap (B \cup C)=&(A\cap B)\cup(A\cap C)\notag   \\
		A\cup (B \cap C)=&(A\cup B)\cap(A\cup C)\notag
	\end{align}}

	除了交并运算，两集合间还有差补运算.
	
	我们定义两集合$S,T$的\textbf{差集}是由属于$S$但不属于$T$的元素组成的集合，记为$S\textbackslash T$，即
	
	\centerline{$S\textbackslash T:=\{ x|x\in S$并且$x\notin T \}$.}
	
	假设我们在集合$X$中讨论，而$S \subset T$，则定义集合$S$关于$X$的\textbf{补集}
	
	\centerline{$\complement_X S:=X\textbackslash S$.}
	
	根据补集的定义，显然有
	
	\centerline{$S \cup \complement_X S=X,S \cup \complement_X S=\varnothing$.}
	
	对于差补的复合运算，可以验证有
	
	\centerline{$S\textbackslash T=S \cap \complement_X T$.}
	
	除了交并差补运算，我们还可以定义集合的\textbf{幂集}.如果我们把集合$A$中每一个子集都当作一个元素，所有子集构成的新集合$P(A)$称为集合$A$的幂集.
	
	显然，当$|A|=n$时，$|P(A)|=2^n$.
	
	\subsection{映射}
	映射是两集合间的一种对应关系.
	
	\paragraph{映射的定义}~{}
	\\
	\indent
	设$X$和$Y$是两个给定的集合，若按照一定的规则$f$，将集合$X$中的每一个元素$x$都惟一对应集合$Y$中的元素$y$，则称这个法则$f$是从集合$X$到$Y$的一个\textbf{映射}，记为
	
	\begin{equation}
		\begin{split}
			f:&X\rightarrow Y\notag   \\
			&x\mapsto y=f(x)\notag 
		\end{split}
	\end{equation}
	
	其中$y$被称为在映射$f$下$x$的\textbf{像}，$x$被称为在映射$f$下$y$的\textbf{原像}，集合$X$称为$f$的\textbf{定义域}，记为$D_f$.在映射$f$下，以集合$X$全体元素$x$的像$y$为元素构成一个新集合，被称为映射$f$的\textbf{值域}，记为$R_f$.
	
	\paragraph{映射的性质}~{}
	\\
	\indent
	按照定义，不难总结出映射三要素：
	
	（1）集合$X$，定义域$D_f=X$；
	
	（2）集合$Y$，值域$R_f \subset Y$;
	
	（3）对应法则$f$，使得任意$x\in X$，有惟一$y=f(x)$与之对应.
	
	映射的定义要求了元素的像具有惟一性，即对任意$x\in X$，对应法则$f$作用后的值惟一.但映射的定义并未对逆像做出要求，逆像并没有惟一性，即允许存在$y\in R_f$，使得$x_1,x_2$都是$y$在$f$下的逆像，$f(x_1)=f(x_2)$.
	
	\paragraph{几种特殊的映射}~{}
	\\
	\indent
	设$f$是集合$X$到集合$Y$的一个映射，若对于$f$的逆像也具有惟一性，则称这样的映射为\textbf{单射}.若该映射$f$的值域$R_f=Y$，则称这样的映射为\textbf{满射}.一个映射若即是单射，又是满射，则称为\textbf{双射}，又称为\textbf{一一对应}.
	
	对于单射$f$，可以构造这样一个映射$g$，使得任意$y\in R_f\subset Y$，有$g(y)=x$.由于单射的定义，这样的$g$是允许存在的.我们把这个从$R_f$到$X$的映射，称作映射$f$的\textbf{逆映射}$f^{-1}$.
	
	可以看出，逆映射的定义域$D_{f^{-1}}=R_f \subset Y$，值域$R_{f^{-1}}=D_f$，则$f^{-1}$是满射，又因为只有单射才能定义逆映射，则逆映射也是单射，则逆映射是$R_f$到$X$的双射.
	
	下面我们来研究映射的乘法（或称为\textbf{复合映射}）.假设有从集合$A$到集合$B$的映射$g$，从集合$C$到集合$D$的映射$f$，若满足$R_g \subset D_f=C$，则可定义这样一个从$A$到$D$的映射，满足元素$x$的像为$f(g(x))$，则称这个映射为映射$f$和映射$g$的乘积（或复合映射），记作$f\circ g(x)$.则有$f\circ g(x)=f(g(x))$.
	
	\paragraph{一元实函数的定义}~{}
	\\
	\indent
	函数是特殊的映射.若取$D_f \subset \mathbb{R}$，$R_f = \mathbb{R}$，则称这样的映射为\textbf{一元实函数}，简称\textbf{函数}.此时的$x$被称为\textbf{自变量}，$y$被称为\textbf{因变量}.
	
	\paragraph{函数的简单特性}~{}
	\\
	\indent
	（1）有界性：若存在两常数$m,M$有，对函数$y=f(x),x\in D$，有$m\le f(x)\ge M,x\in D$，则称函数$f(x)$在$D$上\textbf{有界}.$m$是函数的一个\textbf{下界}，$M$是函数的一个\textbf{上界}.
	
	注意，函数上下界并不唯一，只要常数$m,M$满足上述关系，这样的$m,M$都称为$f(x)$的上下界.
	
	（2）单调性：对函数$y=f(x),x\in D$，任意$x_1,x_2\in D,x_1<x_2$，若$f(x_1)<f(x_2)(f(x_1)\le f(x_2))$，则称$f(x)$在$D$上\textbf{严格单调递增}（\textbf{单调递增}）；反之，若$f(x_1)>f(x_2)(f(x_1)\ge f(x_2))$，则称$f(x)$在$D$上\textbf{严格单调递减}（\textbf{单调递减}）.
	
	（3）奇偶性：当$D$关于原点对称时，即$x\in D \Leftrightarrow -x\in D$.对任意$x\in D$，有$f(x)=f(-x)$，则称$f(x)$为\textbf{偶函数}；有$f(x)+f(-x)=0$，则称$f(x)$为\textbf{奇函数}.
	
	（4）周期性：存在$T>0$，使任意$x\in D$吗，有$f(x)=f(x+T)$，则称$T$是函数$f(x)$的\textbf{周期}，$f(x)$是$D$上的\textbf{周期函数}.若存在一个$T$的最小值，则称这个$T$是函数$f(x)$的\textbf{最小周期}.
	
	\paragraph{复合函数和反函数}~{}
	\\
	\indent
	若函数$f:D\rightarrow f(D)$是单射，则存在逆映射$f^{-1}$，此时的逆映射被称为\textbf{反函数}.
	
	同理，对于两函数$f,g$，若有$R_g \subset D_f$，则可以构造复合映射$f\circ g$，此时的复合映射被称为\textbf{复合函数}.
	
	\paragraph{初等函数}~{}
	\\
	\indent
	我们将下面5种函数叫做\textbf{基本初等函数}:
	
	（1）幂函数：$y=x^\mu(\mu \in \mathbb{R})$，
	
	（2）指数函数：$y=a^x(a>0 $且$ a \neq 1) $，
	
	（3）对数函数：$y=\log_a x(a>0 $且$ a \neq 1)$，
	
	（4）三角函数：$y=\sin x,y=\cos x,y=\tan x,y=\cot x,y=\sec x,y=\csc x$，
	
	（5）反三角函数：$y=\arcsin x,y=\arccos x,y=\arctan x,y=$arccot$\, x,y=$arcsec$\, x,y=$arccsc$\, x$.
	
	由常数和基本初等函数进行有限次四则运算和有限次函数复合运算后的函数叫做\textbf{初等函数}.
	
	
	\section{极限的定义}
	极限的严谨定义是建立在$\epsilon-\delta$语言上的，不论是数列极限还是函数极限，都使用$\epsilon$语言来定义.
	
	\subsection{$\epsilon$语言简述}
	
	下面我们先不加解释的给出这样的符号定义.
	
	$\lim\limits_{x \to x_0}=A \Leftrightarrow \forall \epsilon>0,\,\exists\delta>0,\,\forall x(0<|x-x_0|<\delta):|f(x)-A|<\epsilon$.
	
	观察这一样一个符号定义，我们发现$\epsilon$保证了$f(x)$无限接近于$A$，$\delta$保证了$x$无限接近于$x_0$.即我们可以直观地理解，当$x$趋近于$x_0$时，$f(x)$趋近于$A$.这符合我们直观对极限的理解，又严谨的用数学语言表达了出来.
	
	更广而言，前面的$\forall$后的条件是用来约束因变量的趋近行为，$\exists$后的条件是用来约束自变量的趋近行为.
	
	反之，当我们要证明极限时，常常将$\delta$写成$\delta=\delta(\epsilon)$的形式，通过化简或放缩$|f(x)-A|$，使其在形式上暴露出$|x-x_0|$的形式，再将$\delta$用$\epsilon$表示的形式代入，化简使得不等式右边刚好化为$\epsilon$，使命题得证.下面给出一个例题来解释这样的证明思路.
	
	\textbf{例}$\,\,\,$证明$\lim\limits_{x \to4}\sqrt{x}=2$.
	
	\textbf{分析}$\,\,\,$我们先处理$|f(x)-A|$，即|$\sqrt{x}-2|$的部分，使其暴露出$|x-x_0|$，即$|x-2|$
	
	\centerline{$|\sqrt{x}-2|=\left|\dfrac{x-4}{\sqrt{x}+2}\right|\le \left|\dfrac{x-4}{2}\right|$，}
	
	\noindent{由于要让不等式右边为$\epsilon$，即$\dfrac{|x-4|}{2}\le \epsilon$，反解出$|x-4|=\delta=2\epsilon$.由于此处还要保证$x$的取值在定义域内，代入$x\ge 0$，解的$\delta$的另一个范围$\delta \le 4$.综上，$\delta=\min\{4,2\epsilon\}$.下面写证明的过程即将分析过程反过来即可.}
	
	\textbf{证}$\,\,\,$对$\forall \epsilon>0$，存在$\delta=\min\{4,2\epsilon\}$.使得当$|x-4|<\delta$时，$|\sqrt{x}-2|=\left|\dfrac{x-4}{\sqrt{x}+2}\right|\le \dfrac{|x-4|}{2}= \dfrac{\delta}{2}\le \dfrac{2\epsilon}{2}=\epsilon$.则$\lim\limits_{x \to4}\sqrt{x}=2$.
	
	在证明极限时，我们总是将$\exists$后面条件中的$\delta$（或$X$）用$\forall$后面的$\epsilon$（或$G$）表示，而$\epsilon$的限制不仅仅通过$\epsilon$来限制，还通过函数的自然定义域进行约束.
	\subsection{各种极限的定义}
	\paragraph{数列的极限}~{}
	\\
	\indent
	数列只存在当$n \to \infty$时的极限，则将$\exists$后面的条件更改为$N>0$，而最后$n$应该满足$n>N$.
	
	由于数列中$n\in \mathbb{Z}$，显然，$N\in \mathbb{Z}$，则我们通常将反解出来的$N$加上向上取整函数或者加上取整函数后加1.
	
	\paragraph{函数的极限}~{}
	\\
	\indent
	函数的极限定义形式的多样性来源于$x,f(x)$都既可以趋向于一个具体的值，又可以趋向无穷.
	
	哪个值趋向于无穷，则只需将对应量词后面的条件改成$G>0$或$X>0$，而最后$x,f(x)$应该满足$|x|>X$或$|f(x)|>G$.
	
	
	\section{极限的性质与拓展}
	\subsection{收敛极限的性质}
	极限的性质包含惟一性、（局部）有界性、（局部）保序性、夹逼性.
	\paragraph{惟一性}~{}
	\\
	\indent
	不论是数列还是函数，只要极限存在（即数列为收敛数列或函数在某一点收敛），则极限必\textbf{惟一}.
	
	\paragraph{（局部）有界性}~{}
	\\
	\indent
	对于数列而言，一旦数列收敛，则对$\forall \epsilon$，$\exists N$，使得$\forall n>N$，有$|x_n-A|<\epsilon$.这里的$\epsilon$是任意取得的，则不妨令$\epsilon=1$，则$A-1<x_n<A+1$.则通过收敛数列的定义，则可约束从$N$到无穷大上的所有数落在有界的区间内.而$n<N$时，数的个数是有限的，总能找到一个区间包含这有限的数，再将这两个区间合在一起，则得到：
	
	\centerline{\textbf{收敛数列必有界.}}
	
	对于函数的极限，由于函数研究的是在某一点收敛，则函数在该点的收敛无法推出函数全局有界这么强的结论，但类似存在\textbf{局部有界性}.即对于$\lim\limits_{x \to x_0}f(x)=A$，存在关于$x_0$半径为$\delta$的去心邻域，使得在该去心领域内函数值有界.
	
	\paragraph{（局部）保序性}~{}
	\\
	\indent
	对于数列而言，若$\lim\limits_{n \to \infty}x_n=a,\lim\limits_{n \to \infty}y_n=b$.若$a<b$，则存在$N>0$，使得$n>N$时，$x_n<y_n$.
	
	保序性不难推出保号性，只需将其中一个数列换为全0数列即可.
	
	对于函数而言，若$\lim\limits_{x \to x_0}f(x)=a,\lim\limits_{x \to x_0}g(x)=b$.若$a<b$，则是存在$\delta>0$，使以$\delta$为半径的去心领域内，$f(x)< g(x)$.
	
	\paragraph{夹逼性}~{}
	\\
	\indent
	夹逼性研究的是对极限的性质，因此对数列和函数的极限均适用.
	
	本质上，则是反用极限的保序性，来保证中间这个数列或函数的极限与上下数列或函数的极限相等.
	
	\subsection{收敛数列的判据}
	对于收敛数列的证明，常用的有两种.
	
	\paragraph{单调有界必收敛}~{}
	\\
	\indent
	单调有界数列必定收敛.如果减少一个单调条件，即有界数列，则推导出必有收敛子列.
	
	\paragraph{Cauchy收敛原理}~{}
	\\
	\indent
	若数列$\{x_n\}$满足：$\forall \epsilon>0,\exists N \in \mathbb{N^+}$，使得当$n,m>N$时，有$|x_n-x_m|<\epsilon$.则称数列$\{x_n\}$是\textbf{基本数列}
	
	数列$\{x_n\}$收敛的充要条件是$\{x_n\}$是基本数列.
	
	\subsection{极限的拓展应用}
	\paragraph{无穷小量与无穷大量}~{}
	\\
	\indent
	若$\lim\limits_{x \to x_0}f(x)=0$，则称当$x\to x_0$时，$f(x)$是\textbf{无穷小量}.
	
	高阶无穷小与低阶无穷小则是通过判断两函数趋近于0速度的快慢来定义的，通常通过$\lim\limits_{x \to x_0}\dfrac{f(x)}{g(x)}$的值来确定.
	
	若$\lim\limits_{x \to x_0}f(x)=\infty$，则称当$x\to x_0$时，$f(x)$是\textbf{无穷大量}.
	
	无穷大量的比较与无穷小量一样.
	
	\paragraph{连续函数}~{}
	\\
	\indent
	若函数$f(x)$在点$x_0$邻域上有定义，且成立
	
	\centerline{$\lim\limits_{x \to x_0}f(x)=f(x_0)$}
	
	\noindent{则称$f(x)$在$x_0$\textbf{连续}.}
	
	若$f(x)$在$(a,b)$每一点都连续，则$f(x)$在区间$(a,b)$上连续.
	
	如果仿照左右极限，给出左右连续的定义.$f(x)$在区间$(a,b)$上连续，且$f(x)$在左端点$a$右连续，右端点$a$左连续，则$f(x)$在区间$[a,b]$上连续.
	
	初等函数在其定义域上均为连续函数.
	
	\section{极限的计算}
	\subsection{初等函数代值法}
	在函数的连续中，我们证明了，初等函数都是连续函数，则只要$x$趋近的值在函数定义域内，则直接代值计算.
	
	\textbf{例}$\,\,\,$求$\lim\limits_{x \to 3}\dfrac{x-1}{x+1}$.
	
	
\end{document}